\documentclass{article}

\title{HW7}
\author{Thomas Kidane}
\date{\today}
\usepackage{amsmath,amssymb,graphicx,color,enumitem,tikz}
\begin{document}

\maketitle

\subsection*{Definition 5.5.6}
The random variables $X$ and $Y$, defined on the same sample space, are said to be independent if for all $x$ and $y$,
\[
P(X \leq x, Y \leq y) = P(X \leq x)P(Y \leq y).
\]

\subsection*{Lemma 5.5.9}
The continuous random variables $X$ and $Y$ are independent if and only if they have a joint density $f(x, y)$ which can be written as a product $f(x, y) = g(x)h(y)$ of a function of $x$ alone and a function of $y$ alone. (We say, as before, that $f$ factorises.)

\subsection*{Exercise 5.5.10}
Prove this lemma. You may want to look at the proof of Theorem 2.4.8.

Claim 1($\rightarrow$):
If $f(x,y)$ can be described as a function of $x$ and $y$, then we can get set the distribution function of the two variables as such.

\begin{align*}
  \int_{-\infty}^{\infty}\int_{-\infty}^{\infty}f(x,y)dxdy=1 \rightarrow \int_{-\infty}^{\infty}g(x)dx=1, \int_{-\infty}^{\infty}h(y)dy=1\\
  P(X\leq x)= \int_{-\infty}^{x}g(l)\int_{-\infty}^{\infty}h(y)dydl=\int_{-\infty}^{x}g(l)dl\int_{-\infty}^{\infty}h(y)dy=\int_{-\infty}^{x}g(x)dx\\
  P(Y\leq y)= \int_{-\infty}^{y}h(l)\int_{-\infty}^{\infty}g(x)dxdl=\int_{-\infty}^{y}h(l)dl\int_{-\infty}^{\infty}g(x)dx=\int_{-\infty}^{y}h(y)dy
\end{align*}

But $P(X\leq x, Y\leq y)=\int_{-\infty}^{x}\int_{-\infty}^{y}f(x,y)dxdy=\int_{-\infty}^{x}g(l)dl\int_{-\infty}^{y}h(k)dk=P(X\leq x)P(Y\leq y)$, 
therefore $X$ and $Y$ are independent.

Claim 2($\leftarrow$):

If $X$ and $Y$ are independent then $P(X\leq x,Y\leq y)=P(X\leq x)P(Y\leq y)$. So 
we can make a function $g^{'}(x)=P(X\leq x) \rightarrow g(x)=\frac{d}{dx}P(X\leq x)$, and same definition for $h(y)$.



\end{document}
